% !TeX program = xelatex
\documentclass[10pt]{article}

\newcommand{\ResumeItemSep}{0.05em}
\newcommand{\ResumeTopSep}{0.21em}
\newcommand{\ResumeArrayStretch}{1.06}
\newcommand{\ResumeLineSpread}{1.05}
\newcommand{\ResumeSectionBefore}{0.50em}
\newcommand{\ResumeSectionAfter}{0.20em}

\usepackage{hyperref}
\usepackage{xcolor}
\usepackage{calc}
\usepackage{graphicx}
\usepackage{tikz}
\usepackage{fontspec}
\usepackage{fontawesome5}
\usepackage{titlesec}
\usepackage{enumitem}
\usepackage{fancybox}
\usepackage{xeCJK}

\hypersetup{hidelinks}
\urlstyle{same}

%%%%%%%%%%%%%%%%%%%%
% 设置
%%%%%%%%%%%%%%%%%%%%

\providecommand{\ResumeItemSep}{0.12em}
\providecommand{\ResumeTopSep}{0.18em}
\providecommand{\ResumeArrayStretch}{1.12}
\providecommand{\ResumeLineSpread}{1.08}
\providecommand{\ResumeSectionBefore}{0.55em}
\providecommand{\ResumeSectionAfter}{0.28em}

\setlength{\parindent}{0pt}
\setlength{\parskip}{0pt}
\pagenumbering{gobble}
\setlist[itemize]{
    leftmargin=1.35em,
    itemsep=\ResumeItemSep,
    topsep=\ResumeTopSep,
    parsep=0pt,
    partopsep=0pt
}
\setlist[enumerate]{leftmargin=*}
\renewcommand{\arraystretch}{\ResumeArrayStretch}
\linespread{\ResumeLineSpread}

\titleformat{\section}
  {\Large\bfseries\raggedright}
  {}{0em}
  {}
  [{\color{secondarycolor}\titlerule}]
\titlespacing*{\section}{0pt}{\ResumeSectionBefore}{\ResumeSectionAfter}

\usepackage[
    a4paper,
    left=1.1cm,
    right=1.1cm,
    top=0.75cm,
    bottom=0.85cm,
    nohead
]{geometry}

% 字体设置
% 说明：Noto Serif SC 没有原生斜体字型。将 CJK 的 italic 字型映射到常规字形
%（与过去“找不到斜体、默认替换为常规字形”的实际渲染一致），可消除
% “Font shape `TU/NotoSerifSC(0)/m/it' undefined” 告警；如需真正斜体观感，
% 可在此行追加 ItalicFeatures={FakeSlant=0.18}。
\setCJKmainfont[
    Path=fonts/,
    Extension=.otf,
    BoldFont=*-Bold,
    ItalicFont={NotoSerifSC},
]{NotoSerifSC}

% 自定义颜色
\definecolor{primarycolor}{RGB}{200, 60, 60}
\definecolor{secondarycolor}{RGB}{200, 68, 28}

\newlength{\iconwidth}
\setlength{\iconwidth}{1.5em}

\newcommand{\sectionicon}[2]{\makebox[\iconwidth][c]{\color{primarycolor}{#1}}\quad #2}

\newcommand{\projectheading}[3]{%
    \noindent\makebox[\textwidth][s]{%
        \makebox[0pt][l]{{\large \textbf{#1}}}%
        \hfill
        \makebox[0pt][c]{{\normalsize \textbf{#2}}}%
        \hfill
        \makebox[0pt][r]{#3}%
    }%
}

% 通用个人信息：修改这里即可同步到所有岗位模板。
\newcommand{\ResumeName}{无绝}
\newcommand{\ResumeBirthDate}{2002.10}
\newcommand{\ResumeHometown}{湖北咸宁}
\newcommand{\ResumeSchool}{重庆邮电大学}
\newcommand{\ResumeEnglishLevel}{CET-6}
\newcommand{\ResumeEmail}{250xxxxxxx@qq.com}
\newcommand{\ResumePhone}{198xxxxxxxx}
\newcommand{\ResumeHomepageLabel}{github.com/Dithob}
\newcommand{\ResumeHomepageUrl}{https://github.com/Dithob}
\newcommand{\ResumePhoto}{images/xjpic.jpg}
\newcommand{\ResumeHeaderImage}{images/foot.png}

% 通用教育背景：修改这里即可同步到所有岗位模板。
\newcommand{\ResumeGraduateSchool}{重庆邮电大学}
\newcommand{\ResumeGraduateMajor}{计算机技术}
\newcommand{\ResumeGraduateDegree}{硕士}
\newcommand{\ResumeGraduateDate}{2024.09--2027.06}
\newcommand{\ResumeGraduateAwards}{特等学业奖学金 $\times$ 1、二等学业奖学金 $\times$ 1}

\newcommand{\ResumeUndergraduateSchool}{重庆邮电大学}
\newcommand{\ResumeUndergraduateMajor}{计算机科学与技术}
\newcommand{\ResumeUndergraduateDegree}{本科}
\newcommand{\ResumeUndergraduateDate}{2020.09--2024.07}
\newcommand{\ResumeUndergraduateAwards}{三等学业奖学金；GPA: 3.28/4.0（专业前 25\%）}

\newcommand{\ResumeEducationBlock}{%
    {\textbf{\ResumeGraduateSchool}}，\ResumeGraduateMajor，\textit{\ResumeGraduateDegree} \hfill \ResumeGraduateDate
    \begin{itemize}
        \item \textbf{荣誉奖项}：\ResumeGraduateAwards
    \end{itemize}

    \vspace{0.2em}
    {\textbf{\ResumeUndergraduateSchool}}，\ResumeUndergraduateMajor，\textit{\ResumeUndergraduateDegree} \hfill \ResumeUndergraduateDate
    \begin{itemize}
        \item \textbf{荣誉奖项}：\ResumeUndergraduateAwards
    \end{itemize}
}


%%%%%%%%%%%%%%%%%%%%
% 文章内容
%%%%%%%%%%%%%%%%%%%%

\newcommand{\ResumeTargetRole}{Java 后端开发}
\providecommand{\ResumeTargetRole}{求职方向待定}

\newcommand{\ResumeContact}{%
    \footnotesize
    \textcolor{white}{%
        \faEnvelope \quad 邮箱：\href{mailto:\ResumeEmail}{\ResumeEmail}
        \hspace{3.2em}
        \faPhone \quad 电话：\ResumePhone
        \hspace{3.2em}
        \faGithub \quad 主页：\href{\ResumeHomepageUrl}{\ResumeHomepageLabel}
    }%
}

\newcommand{\ResumeContactBar}{%
    \begin{tikzpicture}[remember picture, overlay]
        \node[anchor=north, inner sep=0pt](headercontacts) at (current page.north){%
            \includegraphics[width=\paperwidth]{\ResumeHeaderImage}
        };
        \node[anchor=center] at(headercontacts.center){\ResumeContact};
    \end{tikzpicture}
    \vspace{-1.15em}
}

\newcommand{\ResumeProfileSection}{%
    \section[个人信息]{\sectionicon{\faAddressCard}{个人信息}}
    \begin{tabular*}{\textwidth}{@{\extracolsep{\fill}}lll}
        \textbf{姓\qquad 名}：\ResumeName & \textbf{出生年月}：\ResumeBirthDate & \textbf{籍\qquad 贯}：\ResumeHometown \\
        \textbf{学\qquad 校}：\ResumeSchool & \textbf{英语水平}：\ResumeEnglishLevel & \textbf{求职方向}：\ResumeTargetRole
    \end{tabular*}
}

\newcommand{\ResumeEducationSection}{%
    \section[教育背景]{\sectionicon{\faGraduationCap}{教育背景}}
    \ResumeEducationBlock
}

\newcommand{\ResumeProfileAndEducation}{%
    %%%%%%%%%%%%%%%%%%%%
    % 个人信息与教育背景
    %%%%%%%%%%%%%%%%%%%%

    \begin{minipage}[t]{0.79\textwidth}
        \ResumeProfileSection

        \ResumeEducationSection
    \end{minipage}
    \hfill
    \begin{minipage}[t]{0.14\textwidth}
        \vspace{-0.35em}
        \setlength{\fboxsep}{0pt}
        % 双线框 (\doublebox) 会给照片额外加约 4pt 边框宽；图片若占满 \linewidth 会触发
        % “Overfull \hbox (3.99988pt too wide)” 告警，故图片宽度预留 4pt，使含框总宽恰好等于栏宽。
        \doublebox{\includegraphics[width=\dimexpr\linewidth-4pt\relax]{\ResumePhoto}}
    \end{minipage}
}


\begin{document}

    % 顶部联系人栏
    \ResumeContactBar
    \ResumeProfileAndEducation

    %%%%%%%%%%%%%%%%%%%%
    % 专业技能
    %%%%%%%%%%%%%%%%%%%%

    \section[专业技能]{\makebox[\iconwidth][c]{\color{primarycolor}{\faWrench}}\quad 专业技能}
    \begin{itemize}
        \item \textbf{Java 与后端框架}：熟悉 Java 集合、异常处理、反射、注解、Stream API 与多线程基础；熟悉 Spring Boot、Spring MVC、MyBatis/MyBatis-Plus，能够完成 RESTful API 设计、参数校验、统一异常处理和业务模块开发。
        \item \textbf{数据库与缓存}：熟悉 MySQL 表结构设计、索引、事务、分页查询和慢 SQL 排查；熟悉 Redis 常用数据结构、缓存穿透/击穿/雪崩处理、Lua 脚本和 Redisson 分布式锁，能够围绕热点数据设计缓存方案。
        \item \textbf{微服务与中间件}：了解 Spring Cloud、Nacos、OpenFeign、Gateway、Sentinel、RabbitMQ、Seata 等组件；具备服务注册发现、统一网关、服务调用、限流降级、异步消息、延迟队列和分布式事务实践。
        % \item \textbf{工程化与质量保障}：熟悉 Linux、Git、Maven、Docker、Postman、JMeter 等工具；能够进行接口联调、日志排查、性能压测、容器化部署和基础运维，关注接口耗时、吞吐量、错误率和数据一致性。
        \item \textbf{AI 应用协作}：具备 RAG 项目实践，熟悉 Linux、Git、Postman、JMeter 等工具；熟悉FastAPI 服务封装、Prompt/RAG 流程调试和模型接口联调，可将 AI 能力以后端 API 形式接入业务系统。
    \end{itemize}

    %%%%%%%%%%%%%%%%%%%%
    % 实习经历
    %%%%%%%%%%%%%%%%%%%%

    \section[实习经历]{\makebox[\iconwidth][c]{\color{primarycolor}{\faBriefcase}}\quad 实习经历}

    \projectheading{腾讯云与智慧产业研发公司}{后端开发实习生}{2026.06--2026.09}
    \begin{itemize}
        \item \textbf{服务开发与工具链}：自研设备控制、线上诊断等 3 个 MCP 服务，接入地图 MCP 动态解析坐标构造测试数据；以 scrcpy 长连接投屏将单次截图耗时从约 1s 降至 50ms。
        \item \textbf{接口质量与缺陷归因}：为 LBS 营销专家（4 个 Skill）搭建 Skill 层 × 接口层双维度评测口径，执行 38 条 P0 用例定位接口层缺陷；完成端到端验收，识别硬编码 Key 泄露等 5 项风险并分级归档。
        \item \textbf{效能与稳定性优化}：参与 Web 自动化框架建设，全量回归由 58.8min 降至 43.6min、执行异常由 34 归零；完成【你用我赔 H5】发版测试，建立跨天状态机、数据一致性等 R1--R8 风险清单。
    \end{itemize}

    % \vspace{0.25em}

    % \projectheading{重庆啄木鸟网络科技有限公司}{后端开发实习生}{2025.06--2025.09}
    % \begin{itemize}
    %     \item \textbf{核心职责}：参与家电售后场景 AI 应用研发，负责需求分析、数据清洗、模型服务接口设计与测试、Prompt/RAG 流程调试、FastAPI 服务封装及前后端联调，支撑智能客服与知识库建设。
    %     \item \textbf{小啄 GPT 智能回复模型}：围绕多轮客服对话中的上下文丢失、槽位遗漏和知识检索不准问题，设计“短期记忆 + 长期摘要 + 结构化状态”机制，用 JSON 维护家电类型、故障现象、预约时间等关键字段；结合 Dense/BM25 双路召回、RRF 融合与 Rerank 精排构建维修知识检索链路，长尾 Query 检索 Hit Rate 提升约 20\%。
    %     \item \textbf{智能产品识别与自动入库}：基于 OCR/视觉大模型、Dify 工作流和向量相似度检索，搭建“图片识别--联网补全--字段校验--重复判断--知识入库”流程；使用 Docker/Miniconda 统一运行环境，并通过 FastAPI 封装识别、补全、去重和入库接口，降低知识库人工维护成本。
    % \end{itemize}

    %%%%%%%%%%%%%%%%%%%%
    % 项目经历
    %%%%%%%%%%%%%%%%%%%%

    \section[项目经历]{\makebox[\iconwidth][c]{\color{primarycolor}{\faProjectDiagram}}\quad 项目经历}

    \projectheading{小麦购票}{后端开发}{2025.09--2025.12}
    \begin{itemize}
        \item \textbf{项目目标}：面向演出购票高并发场景，基于 Spring Cloud、MyBatis-Plus、Nacos、Gateway、Redis、RabbitMQ、Redisson、Caffeine 与 Lua 构建类似大麦的购票系统，覆盖节目展示、票档库存、下单支付和订单超时取消。
        \item \textbf{缓存与热点访问优化}：采用布隆过滤器缓解新用户注册缓存穿透，结合“先查询，后加锁”处理缓存击穿；使用 Caffeine 缓存平台主页数据，将主页访问耗时由 160ms 优化至 90ms；使用 Redis 缓存节目详情和票档余量，将详情查看耗时由 205ms 优化至 132ms，吞吐量由 487/sec 提升至 1493/sec。
        \item \textbf{库存扣减与并发控制}：先使用 Redisson 分布式锁按“节目 + 票档”维度控制购票并发，将吞吐量由 60.1/sec 提升至 68.2/sec；进一步改为 Redis + Lua 原子扣减票档库存，减少锁竞争后吞吐量提升至 88.5/sec。
        \item \textbf{订单可靠性设计}：基于 RabbitMQ 延迟队列实现订单超时自动取消，结合 Saga 思想设计补偿流程，处理支付失败、库存回滚和订单状态修正，降低高并发购票链路中的数据不一致风险。
    \end{itemize}

    \vspace{0.35em}
    \projectheading{基于 RAG 的医药数据分析问数 Agent}{项目主要成员}{2025.10--2026.02}
    \begin{itemize}
        \item \textbf{项目目标}：面向药品采购、零售、库存等全链路数据及 800 余张业务表，基于 Dify、LangChain 与 Milvus 建设 LLM 数据分析问数助手，将业务人员的自然语言问题转化为可执行 SQL 与可读分析结论。
        \item \textbf{问数链路设计}：完成 LLM 应用的工程化部署，设计“问题理解--Schema 检索--SQL 生成--结果解释--人工反馈”端到端链路，打通对话入口、检索服务与 SQL 执行反馈各环节，降低业务人员直接理解复杂数据库结构的门槛。
        \item \textbf{表结构与字段召回}：对数据字典进行元数据增强，将表名、字段名、字段含义和常见枚举值拼接为富文本后向量化，为 Schema 检索提供高质量索引；结合业务规则过滤、问题类型判断、样例对与 CoT，引导模型优先选择相关表和字段，提高复杂多表结构下的 SQL 生成稳定性。
        \item \textbf{执行可靠性与自纠错}：引入 SQL 预执行与自纠错机制，在结果输出前先执行校验，捕获 MySQL Error 后回传模型重写，减少错误结果输出与人工介入，提高复杂多表链路下的执行可靠性。
        \item \textbf{评估体系与指标提升}：构建约 200 条 Golden Dataset，结合 RAGAS 指标、人工反馈 Bad Case 回流，从检索质量与 SQL 生成效果两个维度评估；基于评估结论迭代提示词与检索策略，执行成功率由约 70\% 提升至 90\%+，项目获重庆市人工智能大赛市级银奖。
    \end{itemize}

    % \vspace{0.22em}
    % \projectheading{天天购物平台}{后端开发}{2025.06--2025.08}
    % \begin{itemize}
    %     \item \textbf{项目目标}：基于 Spring Cloud、Nacos、Spring Cloud Gateway、OpenFeign、RabbitMQ、Sentinel、Seata、Redis 与 Redisson 构建类似天猫的电商微服务系统，包含商品展示、搜索、购物车、下单和支付等核心功能。
    %     \item \textbf{微服务架构与治理}：使用 Nacos 完成服务注册发现与配置管理，基于 OpenFeign 实现服务间调用，通过 Gateway 统一路由和请求过滤；接入 Sentinel 完成流量控制与熔断降级，提升高并发场景下服务稳定性。
    %     \item \textbf{性能与异步解耦}：使用 Redis 缓存购物车和热点商品数据，将购物车访问耗时由 119ms 优化至 28ms，商品详情访问耗时由 86ms 优化至 17ms；使用 RabbitMQ 解耦下单/支付流程，通过补偿机制将用户感知支付时间由 3--5s 降至 0.5s 内。
    %     \item \textbf{一致性与部署}：使用 Redisson 分布式锁处理库存扣减，防止高并发超卖；使用 Seata AT 模式处理跨服务分布式事务，并通过 Docker 完成服务容器化部署，降低环境差异带来的发布风险。
    % \end{itemize}
    % ※ JD 切换库存：投电商 / 交易 / 微服务治理方向时，与「医药问数」互换启用。

    % \vspace{0.22em}
    % \projectheading{实时订单系统}{核心开发}{2024.01--2024.04}
    % \begin{itemize}
    %     \item \textbf{项目目标}：基于 Spring Boot、Spring MVC、MyBatis、MySQL、Redis、WebSocket、Spring Task 与微信小程序构建外卖订单系统，覆盖用户点餐、商家管理、订单提醒和运营数据统计。
    %     \item \textbf{核心功能开发}：负责后端架构设计和核心模块开发，实现权限管理、菜品分类、订单状态流转、分页查询和异常处理；使用 Session/Cookie 实现自动登录，使用 Redis 缓存热门菜品信息，降低重复数据库查询。
    %     \item \textbf{实时通知与运营报表}：使用 WebSocket 实现订单提醒和催单功能，提升商家处理订单的实时性；主导数据统计模块开发，按营业额、用户增长和订单完成率等维度提供运营报表，为商家经营分析提供数据支撑。
    % \end{itemize}

    %%%%%%%%%%%%%%%%%%%%
    % 奖项证书
    %%%%%%%%%%%%%%%%%%%%

    \section[在校成果]{\makebox[\iconwidth][c]{\color{primarycolor}{\faAward}}\quad 在校成果}
    \begin{itemize}
        \item \textbf{竞赛}：2025 “互联网+”大学生创新创业大赛银奖；2025 首届重庆市 AI 大模型创新应用大赛省级三等奖；2023 第十四届蓝桥杯 C/C++ 组省赛三等奖；2023 重庆市高校数据库应用竞赛校级二等奖
        \item \textbf{论文}：《LHAF: A Lightweight Hierarchical Adaptive Framework for LLM-based Multimodal Emotion Recognition in Conversations》（SCI I 区在投）
    \end{itemize}

\end{document}
